\section{Conclusiones}
En la configuración evaluada, UMAP retiene más vecinos originales en la AUC que prioriza escalas locales. La ventaja observada se reduce al ponderar las escalas uniformemente y puede invertirse en una escala individual. PCA e ISOMAP tienen una retención media próxima y generan configuraciones geométricas diferentes.

La conservación de vecindarios depende de la escala. La coincidencia semántica es un criterio distinto. UMAP también obtuvo el cambio puntual de pureza ajustada más favorable en la configuración principal, con bandas de diferencias semánticas que incluyen cero. La descomposición exacta atribuye contablemente el cambio de pureza a cuántos vecinos se sustituyen y a la diferencia de pureza entre quienes entran y salen. Las conclusiones describen los grupos léxicos observados y la sensibilidad de estos resultados ante la reponderación de sus documentos.
